\chapter{Sets, Choice and the Real Numbers}

\begin{center}\itshape
Are there more rational numbers than integers? More real numbers than rationals? More points in a square
than on one of its sides? And, before counting them: what \emph{is} a real number?
\end{center}

On 24 November 1858, a young professor at the Zurich Polytechnic was annoyed. Richard Dedekind was teaching
differential calculus for the first time, and he had just told his students that an increasing sequence bounded
above must converge. They believed him; he did not. Every proof he knew leaned on a picture of the line, and a
picture is not a proof. That day he found what was missing, and he waited fourteen years before publishing it, in a
little booklet called \emph{Continuity and Irrational Numbers} (1872).

The same years saw a stranger adventure begin. In November 1873, Georg Cantor wrote to Dedekind with a question
that sounds almost childish: can the natural numbers be paired off with the real numbers? Dedekind answered that
he did not know, but that the algebraic numbers, at least, could be listed. A week later Cantor had the answer,
and it was no: there are \emph{more} real numbers than integers. Infinity comes in sizes. Three years after that he
proved that a square has no more points than one of its sides, and wrote to Dedekind, in French, ``je le vois,
mais je ne le crois pas'': I see it, but I do not believe it.

This chapter follows both stories. It builds the numbers from nothing but sets: the natural numbers, then the
integers and the rationals, then the reals, as limits of rational Cauchy sequences. The axiom that Volume~1 took for
granted, ``every non-empty set of reals bounded above has a least upper bound'', becomes a theorem, and a second
theorem says that there is only one complete ordered field, so it does not matter how the reals are built. Then it
counts, with Cantor's diagonal argument, an idea so good that it came back as Russell's paradox, G\"odel's
incompleteness theorem and Turing's halting problem. It ends with the axiom of choice, which is innocent-looking,
indispensable, and responsible for some of the strangest objects in mathematics.

\section{Sets}

We use sets naively, as everyone does, but it is worth knowing which rules are allowed, because careless rules lead
to contradictions.

\begin{think}
Is there a set of all sets? If $V$ were one, consider $R=\{x\in V:x\notin x\}$. Is $R\in R$?
\end{think}

If $R\in R$, then $R$ satisfies the defining condition, so $R\notin R$; if $R\notin R$, then $R$ satisfies it, so $R\in R$.
This is \emph{Russell's paradox} (1901): the rule ``for every property there is the set of all things having it'' is
contradictory. Zermelo's answer (1908) is to form new sets only from old ones: one may select from an existing set
the elements with a property, but not collect ``all things'' with it. The argument above then proves that there is
no set of all sets.

\paragraph{The axioms of Zermelo--Fraenkel.} Everything is a set, and two sets are equal if they have the same
elements (\emph{extensionality}). The following sets exist: the empty set $\varnothing$; the pair $\{a,b\}$ of two sets;
the union $\bigcup A$ of the elements of a set $A$; the power set $\mathcal P(A)$ of all subsets of $A$; for a set $A$
and a property $P$, the subset $\{x\in A:P(x)\}$ (\emph{separation}); for a set $A$ and a rule $F$ assigning a set $F(x)$
to each $x\in A$, the image $\{F(x):x\in A\}$ (\emph{replacement}); and a set containing $\varnothing$ and stable under
$x\mapsto x\cup\{x\}$ (\emph{infinity}). Finally every non-empty set $A$ has an element disjoint from $A$
(\emph{foundation}; it forbids $x\in x$). These axioms are called ZF. The axiom of choice (AC, Section~\ref{sec:choice})
is added when needed, and ZFC denotes ZF with choice. Ordinary mathematics uses ZFC, and this is the only thing the
volume takes without proof.

\paragraph{Pairs, products, functions.} The ordered pair is $(a,b)=\{\{a\},\{a,b\}\}$ (Kuratowski). Its only purpose is
the property: $(a,b)=(c,d)$ implies $a=c$ and $b=d$. Indeed, if $a=b$ the pair is $\{\{a\}\}$, so $\{c\}=\{c,d\}=\{a\}$ and
$c=d=a$; symmetrically if $c=d$. Otherwise $\{a\}$ is the only one-element member of the left side and $\{c\}$ of the
right, so $a=c$, and then $\{a,b\}=\{a,d\}$ with $b\ne a$ forces $b=d$. The product $A\times B$ is the set of pairs
$(a,b)$, a subset of $\mathcal P(\mathcal P(A\cup B))$, so it exists by separation. A \emph{function} $f:A\to B$ is a subset of
$A\times B$ containing exactly one pair $(a,b)$ for each $a\in A$, and $B^A$ is the set of all functions from $A$ to $B$.
An equivalence relation $\sim$ on $A$ has \emph{classes} $[a]=\{x\in A:x\sim a\}$, which partition $A$, and the
\emph{quotient} $A/{\sim}$ is the set of classes. Quotients are how $\Z$, $\Q$ and $\R$ are built below.

\section{The natural numbers}\label{sec:naturals}

Von Neumann's idea is to let each natural number be the set of the previous ones: $0=\varnothing$, $1=\{0\}$,
$2=\{0,1\}$, and in general $S(n)=n\cup\{n\}$, the \emph{successor} of $n$.

\begin{definition}
A set $I$ is \emph{inductive} if $\varnothing\in I$ and $x\cup\{x\}\in I$ for every $x\in I$. The set $\N$ is the intersection of all
inductive sets: by the axiom of infinity there is an inductive set $I_0$, and $\N=\{x\in I_0:x\text{ belongs to every inductive set}\}$.
\end{definition}

$\N$ is itself inductive, and it is contained in every inductive set. This is the principle of induction in its
purest form: \emph{an inductive subset of $\N$ is $\N$}. In other words, if a property holds for $0$ and passes from $n$
to $S(n)$, it holds on all of $\N$.

\begin{theorem}[Peano's axioms]\label{thm:peano}
$0\ne S(n)$ for all $n\in\N$; $S$ is injective on $\N$; and induction holds.
\end{theorem}
\begin{proof}
$n\in S(n)$ while $0$ has no element, so $S(n)\ne0$. For injectivity, call a set $x$ \emph{transitive} if $y\in x$ implies
$y\subseteq x$. Every $n\in\N$ is transitive: the transitive elements of $\N$ form an inductive set, since $\varnothing$ is transitive
and, if $n$ is, then $y\in S(n)$ means $y\in n$ (so $y\subseteq n\subseteq S(n)$) or $y=n$ (so $y\subseteq S(n)$). Next, $n\notin n$ for $n\in\N$: the
set of such $n$ is inductive, because $\varnothing\notin\varnothing$, and if $n\notin n$ but $S(n)\in S(n)$, then either $S(n)=n$, which gives $n\in n$,
or $S(n)\in n$, which by transitivity gives $n\in S(n)\subseteq n$; both contradict $n\notin n$. Now let $S(m)=S(n)$ with $m\ne n$. Then
$m\in S(n)$, so $m\in n$, and likewise $n\in m$; by transitivity $m\in m$, a contradiction.
\end{proof}

Induction proves properties; \emph{recursion} defines sequences. That a definition ``$f(0)=a$, $f(n+1)=g(f(n))$''
really defines a function is a theorem.

\begin{theorem}[recursion]\label{thm:recursion}
Let $A$ be a set, $a\in A$ and $g:A\to A$. There is a unique $f:\N\to A$ with $f(0)=a$ and $f(S(n))=g(f(n))$ for all $n$.
\end{theorem}
\begin{proof}
Call $h\subseteq\N\times A$ \emph{admissible} if $(0,a)\in h$ and $(n,y)\in h$ implies $(S(n),g(y))\in h$. The set $\N\times A$ is admissible,
and so is the intersection $f$ of all admissible sets. We show that the set $D$ of $n$ for which exactly one $y$ satisfies
$(n,y)\in f$ is inductive; then $f$ is a function on $\N$, and the two equations are the admissibility of $f$.

If $(0,b)\in f$ with $b\ne a$, then $f\setminus\{(0,b)\}$ is still admissible ($0$ is not a successor), contradicting the minimality of $f$; so $0\in D$.
Let $n\in D$, with $(n,y)\in f$. Then $(S(n),g(y))\in f$. If also $(S(n),z)\in f$ with $z\ne g(y)$, then $f'=f\setminus\{(S(n),z)\}$ is admissible: it
contains $(0,a)$ since $S(n)\ne0$, and if $(m,w)\in f'$ then $(S(m),g(w))\in f$ and it differs from $(S(n),z)$, because $S(m)=S(n)$ forces $m=n$
(Theorem~\ref{thm:peano}), hence $w=y$ (as $n\in D$) and $g(w)=g(y)\ne z$. This contradicts minimality again, so $S(n)\in D$.

If $f_1,f_2$ both work, the set of $n$ with $f_1(n)=f_2(n)$ is inductive, so $f_1=f_2$.
\end{proof}

Addition is defined by recursion on the second variable: $m+0=m$ and $m+S(n)=S(m+n)$; multiplication by $m\cdot0=0$ and
$m\cdot S(n)=m\cdot n+m$. Writing $1=S(0)$, we get $S(n)=n+1$.

\begin{proposition}\label{prop:semiring}
Addition and multiplication on $\N$ are associative and commutative, multiplication distributes over addition,
$m+p=n+p$ implies $m=n$, and $m+n=0$ implies $m=n=0$.
\end{proposition}
\begin{proof}
Each statement is proved by induction on one variable, using only the recursive definitions.
(a) $0+n=n$: true for $n=0$, and $0+S(n)=S(0+n)=S(n)$. (b) $S(m)+n=S(m+n)$: true for $n=0$, and
$S(m)+S(n)=S(S(m)+n)=S(S(m+n))=S(m+S(n))$. (c) Commutativity, by induction on $n$: $m+0=m=0+m$ by (a), and
$m+S(n)=S(m+n)=S(n+m)=S(n)+m$ by (b). (d) Associativity, by induction on $p$: $(m+n)+S(p)=S((m+n)+p)=S(m+(n+p))=m+(n+S(p))$.
(e) Cancellation, by induction on $p$: $m+S(p)=n+S(p)$ means $S(m+p)=S(n+p)$, so $m+p=n+p$ by injectivity of $S$.
(f) If $n\ne0$, then $n=S(k)$ for some $k$ (the set $\{0\}\cup S(\N)$ is inductive), and $m+n=S(m+k)\ne0$.

For multiplication: (g) distributivity $m(n+p)=mn+mp$ by induction on $p$, since $m(n+S(p))=m\,S(n+p)=m(n+p)+m=mn+mp+m=mn+m\,S(p)$.
(h) $0\cdot n=0$ and $S(m)\,n=mn+n$, by induction on $n$: $S(m)S(n)=S(m)n+S(m)=mn+n+m+1$ and $m\,S(n)+S(n)=mn+m+n+1$.
(i) Commutativity by induction on $n$: $m\cdot0=0=0\cdot m$, and $m\,S(n)=mn+m=nm+m=S(n)\,m$ by (h). (j) Associativity by
induction on $p$: $(mn)S(p)=(mn)p+mn=m(np)+mn=m(np+n)=m(n\,S(p))$, using (g).
\end{proof}

Order $\N$ by $m\le n$ if $n=m+k$ for some $k\in\N$, and $m<n$ if moreover $k\ne0$.

\begin{proposition}\label{prop:wellorder}
$\le$ is a total order on $\N$, compatible with $+$ and with multiplication by any $p$, and every non-empty subset of $\N$
has a least element.
\end{proposition}
\begin{proof}
Reflexivity and transitivity are clear ($k=0$, and $k+k'$). Antisymmetry: $n=m+k$ and $m=n+k'$ give $m=m+(k+k')$, so $k+k'=0$
by cancellation and $k=k'=0$ by Proposition~\ref{prop:semiring}. Totality, by induction on $n$ for fixed $m$: if $n=0$, then $0\le m$.
If $m<n$ then $m<n+1$; if $m=n$ then $m<n+1$; if $n<m$, then $m=n+S(j)=(n+1)+j$, so $n+1\le m$. Compatibility: $n=m+k$
implies $n+p=(m+p)+k$ and $np=mp+kp$.

Let $A\subseteq\N$ have no least element. We show by induction on $n$ that no $k\le n$ lies in $A$. For $n=0$: $k\le0$ forces $k=0$,
and $0\in A$ would be least. If no $k\le n$ is in $A$ and $n+1\in A$, then $n+1$ is least in $A$, since $k<n+1$ means $k\le n$
(write $n+1=k+S(j)$, so $n=k+j$). So $n+1\notin A$, and no $k\le n+1$ lies in $A$. Hence $A=\varnothing$.
\end{proof}

\section{Integers and rationals}

The integers are formal differences $a-b$ of natural numbers, and two differences are equal if $a+d=b+c$.

\begin{proposition}\label{prop:integers}
On $\N\times\N$, let $(a,b)\sim(c,d)$ if $a+d=b+c$. This is an equivalence relation, and $\Z=(\N\times\N)/{\sim}$ with
\[
\begin{gathered}[(a,b)]+[(c,d)]=[(a+c,b+d)],\qq{}[(a,b)]\cdot[(c,d)]=[(ac+bd,ad+bc)]\end{gathered}
\]
is a commutative ring without zero divisors. The map $n\mapsto[(n,0)]$ embeds $\N$, and every integer is $n$ or $-n=[(0,n)]$ for an $n\in\N$.
With the positive integers $\{[(a,b)]:b<a\}$, $\Z$ is an ordered ring.
\end{proposition}
\begin{proof}
Transitivity: $a+d=b+c$ and $c+f=d+e$ give $a+d+c+f=b+c+d+e$, so $a+f=b+e$ by cancellation. The operations do not depend on
the representatives. For the sum this is direct. For the product, changing $(a,b)$ to $(a',b')$ with $a+b'=a'+b$ gives
\[
\begin{aligned}(ac+bd)+(a'd+b'c)&=c(a+b')+d(a'+b)\\&=c(a'+b)+d(a+b')\\&=(a'c+b'd)+(ad+bc),\end{aligned}
\]
which says $[(ac+bd,ad+bc)]=[(a'c+b'd,a'd+b'c)]$; the second factor is handled symmetrically. Each ring axiom is now an identity
between expressions in $\N$, which follows from Proposition~\ref{prop:semiring}; for instance $[(0,0)]$ is zero and $[(b,a)]$ is the
opposite of $[(a,b)]$. By totality, $(a,b)$ has $a=b+n$ or $b=a+n$, so $[(a,b)]$ is $[(n,0)]$ or $[(0,n)]$. If $x$, $y$ are positive,
$x=[(m,0)]$, $y=[(n,0)]$ with $m,n\ne0$, so $x+y$ and $xy=[(mn,0)]$ are positive (as $mn\ne0$ by Proposition~\ref{prop:semiring}(f) and
(h)). Every non-zero integer is positive or has a positive opposite, so a product of non-zero integers is $\pm$ a positive
one, hence non-zero.
\end{proof}

The rationals are built from $\Z$ exactly as Volume~3, Chapter~6 builds the field of fractions of any domain:
$\Q=\{(a,b)\in\Z\times\Z:b\ne0\}/{\sim}$ with $(a,b)\sim(c,d)$ iff $ad=bc$, and the usual operations on $\frac ab$. The verifications are
the same kind of identities as for $\Z$, and the result is a field containing $\Z$. Order it by: $\frac ab>0$ iff $ab>0$, which
does not depend on the representative (if $\frac ab=\frac cd$, then $ab\,d^2=cd\,b^2$, and $b^2,d^2>0$). Positive rationals are
stable under sum and product, and $\Q$ is an \emph{ordered field}. It is \emph{Archimedean}: for $\frac ab>0$ with $a,b>0$, the
integer $n=a+1$ satisfies $n>\frac ab$, since $nb-a\ge a+1-a>0$.

\begin{definition}
An \emph{ordered field} is a field $F$ with a subset $F_{>0}$, stable under sums and products, such that each $x\in F$ is in
exactly one of $F_{>0}$, $\{0\}$, $-F_{>0}$; one writes $x<y$ for $y-x\in F_{>0}$. It is \emph{Archimedean} if for every $x\in F$ there is
$n\in\N$ with $x<n\cdot1$; it has the \emph{least upper bound property} if every non-empty subset bounded above has a least upper
bound.
\end{definition}

In an ordered field $1=1^2>0$, so $n\cdot1>0$ for $n\ge1$: the map $\Z\to F$ is injective, and $F$ contains a copy of $\Q$. The rules
of inequalities of Volume~1 hold in any ordered field, with the same proofs. $\Q$ lacks the least upper bound property:
$\{q\in\Q:q^2<2\}$ has no least upper bound in $\Q$ (Exercise~\ref{ex:sqrt2}).

\section{The real numbers}\label{sec:reals}

How do you build a number that is not there? The rationals have holes: the sequence $1$, $1.4$, $1.41$, $1.414$, \ldots{}
crowds together as if it were going somewhere, and in $\Q$ there is nowhere to go. Dedekind's answer was to define
$\sqrt2$ as the cut it makes in the rationals, the pair (those below, those above). Cantor's answer, which we follow, was
bolder and lazier at the same time: \emph{define} the missing limit to be the sequence itself, or rather the whole
family of sequences that crowd together in the same way.

Cantor (1872) and M\'eray (1869) defined a real number as what a Cauchy sequence of rationals ``converges to'',
identifying two sequences when their difference tends to $0$. In this section, $\varepsilon$ and $\delta$ denote positive
rationals, and sequences are sequences of rationals unless said otherwise.

\begin{definition}
A sequence $(x_n)$ in $\Q$ is \emph{Cauchy} if for every $\varepsilon$ there is $N$ with $|x_n-x_m|<\varepsilon$ for all $n,m\ge N$, and \emph{null}
if for every $\varepsilon$ there is $N$ with $|x_n|<\varepsilon$ for $n\ge N$. Let $\mathcal C$ be the set of Cauchy sequences, $\mathcal N$ that of null
sequences, and $\R=\mathcal C/{\sim}$, where $x\sim y$ iff $x-y\in\mathcal N$. The class of $x$ is written $[x]$.
\end{definition}

\begin{lemma}\label{lem:cauchybasic}
\begin{enumerate}[label=(\alph*)]
\item A Cauchy sequence is bounded.
\item $\mathcal C$ is a commutative ring for termwise operations, and $\mathcal N$ is an ideal of it.
\item If $x\in\mathcal C\setminus\mathcal N$, there are $\delta$ and $N$ such that either $x_n\ge\delta$ for all $n\ge N$, or $x_n\le-\delta$ for all $n\ge N$.
\end{enumerate}
\end{lemma}
\begin{proof}
(a) With $\varepsilon=1$, $|x_n|\le|x_N|+1$ for $n\ge N$, and there are finitely many earlier terms. (b) Sums are clear. For
products, if $|x_n|,|y_n|\le M$ for all $n$, then $|x_ny_n-x_my_m|\le M|y_n-y_m|+M|x_n-x_m|$. The same inequality with $y$ null shows
that a bounded sequence times a null sequence is null, so $\mathcal N$ is an ideal. (c) Since $x$ is not null, there is $\varepsilon$ with
$|x_n|\ge\varepsilon$ for infinitely many $n$. Let $\delta=\frac\varepsilon2$ and $N$ with $|x_n-x_m|<\delta$ for $n,m\ge N$, and pick $m\ge N$ with $|x_m|\ge\varepsilon$. If
$x_m\ge\varepsilon$ then $x_n>x_m-\delta\ge\delta$ for all $n\ge N$; if $x_m\le-\varepsilon$ then $x_n\le-\delta$ similarly.
\end{proof}

So $\R=\mathcal C/\mathcal N$ is a commutative ring, and $\Q$ maps into it by constant sequences, $q\mapsto\iota(q)$. Call $[x]$ \emph{positive}
if the first case of (c) occurs. This does not depend on the representative: adding a null sequence $z$, with
$|z_n|<\frac\delta2$ for large $n$, keeps $x_n+z_n\ge\frac\delta2$.

\begin{theorem}\label{thm:reals}
$\R$ is an Archimedean ordered field, $\iota:\Q\to\R$ is an injective morphism of ordered fields, and $\Q$ is dense: between
two reals $a<b$ there is a rational $q$ with $a<\iota(q)<b$. Every Cauchy sequence of $\R$ converges in $\R$, and $\R$ has the least
upper bound property.
\end{theorem}
\begin{proof}
\emph{Field.} If $[x]\ne0$, Lemma~\ref{lem:cauchybasic}(c) gives $|x_n|\ge\delta$ for $n\ge N$. Let $y_n=\frac1{x_n}$ for $n\ge N$ and $y_n=1$ before.
Then $|y_n-y_m|=\frac{|x_m-x_n|}{|x_nx_m|}\le\frac{|x_m-x_n|}{\delta^2}$ for $n,m\ge N$, so $y\in\mathcal C$, and $x_ny_n-1=0$ for $n\ge N$: $[x][y]=1$. Also
$1\ne0$, since the constant sequence $1$ is not null.

\emph{Order.} By Lemma~\ref{lem:cauchybasic}(c), each non-zero $[x]$ is positive or has a positive opposite, and not both, because
$x_n\ge\delta$ and $x_n\le-\delta'$ cannot both hold for large $n$. Sums and products of positive elements are positive
($x_n+y_n\ge\delta+\delta'$ and $x_ny_n\ge\delta\delta'$). The map $\iota$ is a ring morphism, injective because a non-zero constant is not null, and $q>0$
iff $\iota(q)>0$.

\emph{Archimedes.} If $|x_n|\le M$ for all $n$ (Lemma~\ref{lem:cauchybasic}(a)), then $M+1-x_n\ge1$, so $[x]<\iota(M+1)$, and $M+1<n$ for some $n\in\N$.

\emph{Density.} Let $a=[x]<b=[y]$: there are $\delta$ and $N_0$ with $y_n-x_n\ge\delta$ for $n\ge N_0$. Choose $N\ge N_0$ with
$|x_n-x_m|,|y_n-y_m|<\frac\delta4$ for $n,m\ge N$, and let $q=\frac{x_N+y_N}2$. For $n\ge N$,
$q-x_n\ge q-x_N-\frac\delta4=\frac{y_N-x_N}2-\frac\delta4\ge\frac\delta4$, so $\iota(q)>a$; likewise $\iota(q)<b$.

\emph{Rational approximations converge.} Let $x\in\mathcal C$. Given $\varepsilon$, take $N$ with $|x_n-x_m|\le\frac\varepsilon2$ for $n,m\ge N$. For fixed $n\ge N$,
the sequence $(\frac\varepsilon2-|x_m-x_n|)_m$ is $\ge0$ for $m\ge N$, so $|[x]-\iota(x_n)|\le\iota(\frac\varepsilon2)<\iota(\varepsilon)$. Every real $e>0$ exceeds some
$\iota(\varepsilon)$ by density, so $\iota(x_n)\to[x]$ in $\R$.

\emph{Completeness.} Let $(a_k)$ be Cauchy in $\R$. By density pick rationals $q_k$ with $|a_k-\iota(q_k)|<\iota(\frac1k)$. Then
$|q_k-q_l|\le|q_k-a_k|+|a_k-a_l|+|a_l-q_l|$ shows that $(q_k)$ is Cauchy in $\Q$ (an inequality between rationals holds iff it holds after
$\iota$). Let $a=[(q_k)]$. Then $|a-a_k|\le|a-\iota(q_k)|+\frac1k\to0$ by the previous step.

\emph{Least upper bounds.} Let $A\ne\varnothing$ be bounded above by $u_0$, and let $l_0$ be smaller than some element of $A$. Bisect: with
$m=\frac{l_k+u_k}2$, set $(l_{k+1},u_{k+1})=(l_k,m)$ if $m$ bounds $A$ above, and $(m,u_k)$ otherwise. Then $u_k$ bounds $A$, $l_k$ does not, and
$u_k-l_k=2^{-k}(u_0-l_0)$, which tends to $0$ by Archimedes. Both sequences are monotone with $l_k\le l_j\le u_j\le u_k$ for $j\ge k$, hence Cauchy,
with a common limit $s$. For $a\in A$, $a\le u_k$ for all $k$, so $a\le s$ (if $a>s$, then $u_k<a$ for large $k$). If $v<s$, then $v<l_k$ for some
$k$, and some $a\in A$ exceeds $l_k>v$, so $v$ is not an upper bound. So $s=\sup A$.
\end{proof}

From now on we identify $q$ with $\iota(q)$. The axiom of Volume~1 is now a theorem. It is the only one needed: every
other property of $\R$ used in the series follows from it, and the next theorem shows that it determines $\R$.

\begin{theorem}[uniqueness of $\R$]\label{thm:uniqueR}
Let $F$ be an ordered field with the least upper bound property. There is a unique field isomorphism $\varphi:\R\to F$, and it
preserves the order.
\end{theorem}
\begin{proof}
\emph{$F$ is Archimedean}: otherwise $\N\cdot1$ is bounded above, with supremum $s$; then $s-1<n$ for some $n\in\N$, and $n+1>s$. \emph{$\Q$ is
dense in $F$}: for $x<y$, take $n>\frac1{y-x}$, and let $m$ be the least integer with $m>nx$ (Archimedes and Proposition~\ref{prop:wellorder});
then $nx<m\le nx+1<ny$, so $x<\frac mn<y$. The same holds in $\R$.

For $x\in\R$ let $L(x)=\{q\in\Q:q<x\}$ and $\varphi(x)=\sup_FL(x)$, which exists ($L(x)$ is non-empty and bounded by any rational $\ge x$). For
$x\in\Q$, $\varphi(x)=x$ by density in $F$. If $x<y$, two rationals $x<q<q'<y$ give $\varphi(x)\le q<q'\le\varphi(y)$: $\varphi$ is strictly increasing.
Moreover $\{q\in\Q:q<\varphi(x)\}=L(x)$: a rational $q<x$ has a rational $q<q'<x$, so $q<q'\le\varphi(x)$; and a rational $q\ge x$ bounds
$L(x)$, so $\varphi(x)\le q$. It is onto: for $y\in F$, let $x=\sup_\R\{q\in\Q:q<y\}$; the same argument shows that $L(x)=\{q\in\Q:q<y\}$, so
$\varphi(x)=\sup_FL(x)=y$ by density in $F$.

\emph{Sums.} $L(x+y)=L(x)+L(y)$: if $q<x+y$, pick a rational $r$ with $q-y<r<x$, then $q=r+(q-r)$ with $q-r<y$; the other inclusion is clear.
The supremum of a sumset is the sum of the suprema (in any ordered field: $a+b\le\sup A+\sup B$, and $\sup A-\eta<a$, $\sup B-\eta<b$
gives $a+b>\sup A+\sup B-2\eta$), so $\varphi(x+y)=\varphi(x)+\varphi(y)$, and then $\varphi(-x)=-\varphi(x)$.
\emph{Products.} For $x,y>0$, the positive part of $L(xy)$ is $\{rs:r\in L(x),s\in L(y),r,s>0\}$, by the same argument with quotients, and
suprema of products of positive sets multiply. So $\varphi(xy)=\varphi(x)\varphi(y)$ for $x,y>0$, and for all signs by $\varphi(-x)=-\varphi(x)$.

\emph{Uniqueness.} Let $\psi:\R\to F$ be any field morphism. It fixes $\Z$, hence $\Q$. Every $x\ge0$ in $\R$ is a square: $s=\sup\{q\in\Q:q^2<x\}$
satisfies $s^2=x$ (if $s^2<x$, then $(s+h)^2<x$ for a small rational $h>0$; if $s^2>x$, then $(s-h)^2>x$ for a small $h$; either
contradicts the definition of $s$). So $x\ge0$ gives $\psi(x)=\psi(\sqrt x)^2\ge0$: $\psi$ is increasing. Then for $q<x<q'$ rational,
$q\le\psi(x)\le q'$, so $\psi(x)=\varphi(x)$ by density of $\Q$ in $F$.
\end{proof}

The second road to $\R$, Dedekind's cuts (1872), defines a real number as the set of rationals below it; by this
theorem it gives the same field (Exercise~\ref{ex:cuts}). A third road is the decimal expansions of school; it works
too, but carries are painful.

\section{Countable and uncountable}

Imagine a hotel with infinitely many rooms, numbered $0,1,2,\ldots$, all occupied. A new guest arrives. The manager is not
worried: she asks every guest to move from room $n$ to room $n+1$, and room $0$ is free. A coach with infinitely many new
guests arrives; she moves the guest of room $n$ to room $2n$, and the odd rooms are free. Hilbert liked to tell this story
in his lectures (1924) to show that ``as many as'' is a subtle notion for infinite sets. The right definition is
Cantor's: two sets have the same size if their elements can be paired off.

\begin{definition}
Two sets $A$, $B$ are \emph{equipotent}, $|A|=|B|$, if there is a bijection $A\to B$; and $|A|\le|B|$ if there is an injection $A\to B$. A set
is \emph{finite} if it is equipotent to some $n=\{0,\dots,n-1\}$, and \emph{countable} if $|A|\le|\N|$.
\end{definition}

\begin{think}
Are there more rationals than integers? There are infinitely many rationals between any two integers. Are there
more real numbers than rationals, although every real is a limit of rationals?
\end{think}

\begin{proposition}\label{prop:countable}
\begin{enumerate}[label=(\alph*)]
\item A subset of $\N$ is finite or equipotent to $\N$; so a countable set is finite or in bijection with $\N$.
\item $\N\times\N$ is countable, and so are $\Z$ and $\Q$.
\item (Using countable choice.) A countable union of countable sets is countable.
\end{enumerate}
\end{proposition}
\begin{proof}
(a) Let $A\subseteq\N$ be infinite. By recursion, let $a_0=\min A$ and $a_{k+1}=\min\{a\in A:a>a_k\}$, which exists since $A$ is infinite
(Proposition~\ref{prop:wellorder}). Then $k\mapsto a_k$ is strictly increasing, so $a_k\ge k$, and every $a\in A$ is some $a_k$: otherwise take the least
$k$ with $a_k>a$; then $k\ge1$ (as $a_0\le a$) and $a_{k-1}<a<a_k$ contradicts the definition of $a_k$. (b) $(m,n)\mapsto2^m3^n$ is injective by
unique factorisation. $\Z$ injects into $\N\times\N$ by $n\mapsto(n,0)$, $-n\mapsto(0,n)$, and $\Q$ into $\Z\times\Z$ by writing fractions in lowest terms with a
positive denominator. (c) Let $A_0,A_1,\ldots$ be countable, and \emph{choose} for each $k$ an injection $f_k:A_k\to\N$. For $a\in\bigcup A_k$ let $k(a)$ be the
least $k$ with $a\in A_k$; then $a\mapsto(k(a),f_{k(a)}(a))$ injects $\bigcup A_k$ into $\N\times\N$.
\end{proof}

So $\Q$ can be listed as $q_0,q_1,q_2,\ldots$; the Calkin--Wilf tree of Figure~\ref{fig:calkinwilf} lists the positive rationals
without repetition, and even with a formula for the next term (Exercise~\ref{ex:calkinwilf}).

\begin{figure}[ht]
\centering
\fig{ch01-calkinwilf}
\caption{The Calkin--Wilf tree: the children of $\frac ab$ are $\frac a{a+b}$ and $\frac{a+b}b$. Every positive rational appears exactly once, in lowest
terms, and the path from the root to $\frac ab$ is Euclid's algorithm run backwards.}
\label{fig:calkinwilf}
\end{figure}

\begin{theorem}[Cantor, 1891]\label{thm:cantor}
For every set $A$ there is no surjection $A\to\mathcal P(A)$. In particular $|A|<|\mathcal P(A)|$, and $\mathcal P(\N)$ is uncountable.
\end{theorem}
\begin{proof}
Let $f:A\to\mathcal P(A)$ and $D=\{a\in A:a\notin f(a)\}$. If $D=f(d)$, then $d\in D$ iff $d\notin f(d)=D$, a contradiction. So $D$ is not in the image.
And $a\mapsto\{a\}$ is an injection $A\to\mathcal P(A)$.
\end{proof}

\begin{figure}[ht]
\centering
\fig{ch01-diagonal}
\caption{Cantor's diagonal argument for binary sequences, that is, subsets of $\N$. Given any list $s_0,s_1,\ldots$, flip the diagonal: the
sequence $d$ differs from $s_n$ at place $n$, so it is not in the list. The proof of Theorem~\ref{thm:cantor} is the same argument, with
$a\in f(a)$ in place of the digit.}
\label{fig:diagonal}
\end{figure}

The argument is the same as Russell's, and the same as Turing's proof that no program decides whether programs halt
. To compare $\R$ with $\mathcal P(\N)$ we need a way to build bijections from two injections.

\begin{theorem}[Schr\"oder--Bernstein]\label{thm:sb}
If there are injections $f:A\to B$ and $g:B\to A$, then there is a bijection $A\to B$.
\end{theorem}
\begin{proof}
Let $C_0=A\setminus g(B)$, $C_{n+1}=g(f(C_n))$ and $C=\bigcup_nC_n$. Define $h(a)=f(a)$ if $a\in C$ and $h(a)=g^{-1}(a)$ if $a\notin C$; this makes sense
because $a\notin C$ implies $a\notin C_0$, so $a\in g(B)$. \emph{Injective}: $h$ is injective on $C$ and on $A\setminus C$. If $f(a)=g^{-1}(a')$ with $a\in C_n$ and
$a'\notin C$, then $a'=g(f(a))\in C_{n+1}$, a contradiction. \emph{Surjective}: let $b\in B$. If $g(b)\notin C$, then $h(g(b))=b$. If $g(b)\in C$, then
$g(b)\notin C_0$, so $g(b)\in C_{n+1}=g(f(C_n))$ for some $n$, and since $g$ is injective, $b=f(a)$ with $a\in C_n$; then $h(a)=b$.
\end{proof}

\begin{figure}[ht]
\centering
\fig{ch01-sb}
\caption{The proof of Schr\"oder--Bernstein. The elements of $A$ not hit by $g$ start chains $a\mapsto f(a)\mapsto g(f(a))\mapsto\cdots$; on the elements of these
chains (the set $C$, in coral) the bijection is $f$, everywhere else it is $g^{-1}$.}
\label{fig:sb}
\end{figure}

\begin{theorem}\label{thm:continuum}
$|\R|=|[0,1]|=|\mathcal P(\N)|$. In particular $\R$ is uncountable, and so is every interval with more than one point.
\end{theorem}
\begin{proof}
$x\mapsto L(x)=\{q\in\Q:q<x\}$ is injective from $\R$ to $\mathcal P(\Q)$, by density of $\Q$, and $|\mathcal P(\Q)|=|\mathcal P(\N)|$ since a bijection $\Q\to\N$
(Proposition~\ref{prop:countable}) induces one on subsets. Conversely $A\mapsto\sum_{n\in A}2\cdot3^{-n-1}$ is an injection $\mathcal P(\N)\to[0,1]$: if $A\ne B$ and $n$ is the
least element of their symmetric difference, say $n\in A$, the two sums differ by at least $2\cdot3^{-n-1}-\sum_{k>n}2\cdot3^{-k-1}=3^{-n-1}>0$. So
\[|\mathcal P(\N)|\le|[0,1]|\le|\R|\le|\mathcal P(\N)|,\]
and Schr\"oder--Bernstein gives the equalities. An interval $[a,b]$ with $a<b$ is in bijection with $[0,1]$.
\end{proof}

Every countable subset of $\R$ therefore misses most reals. With the algebraic numbers, which are countable
(Exercise~\ref{ex:algebraic}), this proves that transcendental numbers exist, without exhibiting a single one; Cantor's
1874 paper did exactly this. Is there a set of reals strictly between $\N$ and $\R$ in size? The \emph{continuum
hypothesis} says no. G\"odel (1938) and Cohen (1963) showed that ZFC can neither prove nor refute it.

\section{Choice, Zorn and well-ordering}\label{sec:choice}

In 1904, Ernst Zermelo proved that every set can be well-ordered, every set, $\R$ included. The proof took two pages
and caused a storm. It rested on a principle that mathematicians had used silently for decades: from any family of
non-empty sets one can choose one element in each. Borel, Baire and Lebesgue protested that a ``choice'' nobody can
describe is no choice at all; Hadamard answered that existence and description are different things. A century later
the verdict is in: the axiom of choice cannot be proved or refuted from the other axioms (G\"odel 1938, Cohen 1963),
nearly all of mathematics uses it, and it has consequences that defy intuition. The most famous, the Banach--Tarski
paradox, cuts a ball into five pieces and reassembles them into two balls of the same size. (Question: what is an
anagram of ``Banach--Tarski''? Answer: ``Banach--Tarski Banach--Tarski''.)

\begin{definition}
The \emph{axiom of choice} (AC) says: for every set $\mathcal A$ of non-empty sets there is a function $c$ on $\mathcal A$ with $c(A)\in A$ for all
$A\in\mathcal A$ (a \emph{choice function}).
\end{definition}

For finitely many sets, choice follows from the other axioms by induction; for the countably many enumerations in
Proposition~\ref{prop:countable}(c), or for one point in each class of $\R/\Q$ (Vitali's set, Chapter~7), it does not. AC is most
often used through an equivalent statement about orders.

\begin{definition}
A \emph{partial order} on $P$ is reflexive, antisymmetric and transitive. A \emph{chain} is a totally ordered subset. An \emph{upper bound} of
$C\subseteq P$ is a $u$ with $c\le u$ for all $c\in C$; a \emph{maximal} element is an $m$ with no $x>m$. A \emph{well-order} is a total order in which every
non-empty subset has a least element.
\end{definition}

\begin{theorem}[Zorn's lemma]\label{thm:zorn}
Assume AC. If every chain of a non-empty partially ordered set $P$ has an upper bound, then $P$ has a maximal element.
\end{theorem}
\begin{proof}
Suppose that $P$ has no maximal element. Then every chain $C$ has a \emph{strict} upper bound: an upper bound $u$ exists, $u$ is not
maximal, so some $v>u$ bounds $C$ strictly. By AC choose a strict upper bound $g(C)$ for every chain $C$ (the empty chain included).
Call $W\subseteq P$ a \emph{$g$-set} if it is well-ordered by $\le$ and $w=g(W_{<w})$ for every $w\in W$, where $W_{<w}=\{x\in W:x<w\}$. A subset $S\subseteq W$
is an \emph{initial segment} of $W$ if $s\in S$, $w\in W$ and $w<s$ imply $w\in S$.

\emph{Two $g$-sets $W$, $W'$ are comparable: one is an initial segment of the other.} Let $D$ be the union of all sets that are initial
segments of both; $D$ is one too. Suppose $D\ne W$ and $D\ne W'$, and let $w=\min(W\setminus D)$, $w'=\min(W'\setminus D)$. Then $W_{<w}=D$: elements below $w$
are in $D$ by minimality, and if $d\in D$ had $d\ge w$, then $w\in D$ since $D$ is an initial segment. Likewise $W'_{<w'}=D$, so $w=g(D)=w'$. Then
$D\cup\{w\}$ is an initial segment of both $W$ and $W'$, contradicting the definition of $D$.

\emph{The union $U$ of all $g$-sets is a $g$-set.} Two elements of $U$ lie in $g$-sets $W$, $W'$, one of which contains the other, so they
are comparable: $U$ is totally ordered. Each $g$-set $W$ is an initial segment of $U$: if $w\in W$ and $u<w$ with $u$ in a $g$-set $W'$,
then $u\in W$ whether $W'$ is an initial segment of $W$ or $W$ one of $W'$. Hence a non-empty $S\subseteq U$, meeting some $W$, has as least element
the least element of $S\cap W$; and for $u\in W$, $U_{<u}=W_{<u}$, so $g(U_{<u})=u$.

Now $U$ is a chain, so $g(U)$ is defined and $g(U)\notin U$. But $U\cup\{g(U)\}$ is a $g$-set (with $g(U)$ on top, and $g(U)=g(U_{<g(U)})$), so it is
contained in $U$: a contradiction.
\end{proof}

\begin{theorem}\label{thm:wellordering}
In ZF, the following are equivalent: AC; Zorn's lemma; every set can be well-ordered (Zermelo, 1904).
\end{theorem}
\begin{proof}
AC implies Zorn by Theorem~\ref{thm:zorn}. \emph{Zorn implies well-ordering}: given $X$, let $P$ be the set of pairs $(S,\preceq)$ where $S\subseteq X$
and $\preceq$ well-orders $S$, ordered by $(S,\preceq)\le(S',\preceq')$ if $S$ is an initial segment of $S'$ and $\preceq'$ extends $\preceq$. A chain has as upper
bound its union $S^*$ with the union order: it is a total order, and a non-empty $A\subseteq S^*$ meets some $S$ of the chain, which is an initial
segment of $S^*$, so the least element of $A\cap S$ is least in $A$. Let $(M,\preceq)$ be maximal. If $x\in X\setminus M$, putting $x$ above $M$ gives a larger
element of $P$; so $M=X$. \emph{Well-ordering implies AC}: well-order $\bigcup\mathcal A$ and let $c(A)$ be the least element of $A$.
\end{proof}

\begin{theorem}\label{thm:basis}
Assume AC. Every vector space $V$ over a field $K$ has a basis; more precisely every linearly independent subset of $V$ is
contained in a basis.
\end{theorem}
\begin{proof}
Let $L_0$ be independent, and order the independent sets containing $L_0$ by inclusion. The union of a chain is independent:
a linear relation involves finitely many vectors, which all lie in one member of the chain. By Zorn there is a maximal
independent $B\supseteq L_0$. If $v\notin\operatorname{span}B$, then $B\cup\{v\}$ is independent, contradicting maximality. So $B$ spans $V$.
\end{proof}

A basis of $\R$ as a vector space over $\Q$ is called a \emph{Hamel basis} (1905). Nobody can write one down, and it gives
strange functions (Exercise~\ref{ex:hamel}). The same pattern proves that every proper ideal of a commutative ring is
contained in a maximal ideal (Exercise~\ref{ex:krull}), that every field has an algebraic closure, and, in Chapter~4,
Tychonoff's theorem. AC is independent of ZF: G\"odel (1938) showed that it cannot be refuted, Cohen (1963) that it cannot
be proved.

\begin{beauty}
Numbers, functions and spaces are all sets, and a handful of axioms produces all of them. The construction of $\R$
from $\Q$ is the first \emph{completion}: Chapter~2 repeats it for any metric space and finds the $p$-adic numbers of
Volume~2 as completions of $\Q$ for other distances, and Volume~8 finds $L^2$ as the completion of continuous
functions. Cantor's diagonal is the most portable argument in mathematics: it proves that $\R$ is uncountable, that
there is no set of all sets, that some true statements of arithmetic are unprovable (G\"odel) and that some
problems cannot be solved by any program (Turing), foundational results in mathematical logic. And the axiom of choice shows where
intuition ends: it gives every vector space a basis, and it gives subsets of the line with no reasonable length.
\end{beauty}

\section{Exercises}

\begin{exercise}[$\star$]\label{ex:sqrt2}
Show that $A=\{q\in\Q:q>0,\ q^2<2\}$ has no least upper bound in $\Q$. (Show that $q\mapsto\frac{2q+2}{q+2}$ maps $A$ into $A$ and increases every element,
and maps the rational upper bounds of $A$ to smaller ones.)
\end{exercise}

\begin{exercise}[$\star$]\label{ex:algebraic}
A complex number is \emph{algebraic} if it is a root of a non-zero polynomial with integer coefficients. Show that the
algebraic numbers form a countable set. Deduce that the transcendental real numbers form an uncountable set.
\end{exercise}

\begin{exercise}[$\star$]\label{ex:cauchyfe}
Let $f:\R\to\R$ satisfy $f(x+y)=f(x)+f(y)$ for all $x,y$. Show that $f(q)=qf(1)$ for $q\in\Q$, and that $f(x)=xf(1)$ for all $x$ if $f$ is continuous,
or monotone.
\end{exercise}

\begin{exercise}[$\star\star$]\label{ex:square}
Show that $|[0,1]\times[0,1]|=|[0,1]|$: a square has as many points as one of its sides. Deduce that $|\R^n|=|\R|$ for every $n\ge1$.
\end{exercise}

\begin{exercise}[$\star\star$]\label{ex:hamel}
(a) Show that $\R$ has a basis $B$ as a vector space over $\Q$. (b) Fix $b_0\in B$, and let $f(x)$ be the coordinate of $x$ on $b_0$. Show that $f$ is additive
but not of the form $x\mapsto cx$. (c) Show that the graph of every additive $f:\R\to\R$ that is not of the form $x\mapsto cx$ is dense in $\R^2$. Deduce that
such an $f$ is unbounded on every interval.
\end{exercise}

\begin{exercise}[$\star\star$]\label{ex:cuts}
A \emph{Dedekind cut} is a set $\alpha\subseteq\Q$ that is non-empty, not all of $\Q$, closed downwards ($q<a\in\alpha$ implies $q\in\alpha$) and without a largest
element. (a) Show that the set of cuts, ordered by inclusion, has the least upper bound property. (b) Show that $x\mapsto\{q\in\Q:q<x\}$ is an
increasing bijection from $\R$ onto the set of cuts.
\end{exercise}

\begin{exercise}[$\star\star$]\label{ex:krull}
Assume AC. Let $R$ be a commutative ring with $1\ne0$. Show that every ideal $I\ne R$ is contained in a maximal ideal.
\end{exercise}

\begin{exercise}[$\star\star$]\label{ex:functions}
Show that the set $C(\R,\R)$ of continuous functions has the cardinality of $\R$, while the set $\R^\R$ of all functions has the
cardinality of $\mathcal P(\R)$, which is larger.
\end{exercise}

\begin{exercise}[$\star\star\star$]\label{ex:calkinwilf}
\emph{Recounting the rationals.} In the Calkin--Wilf tree (Figure~\ref{fig:calkinwilf}), the root is $\frac11$ and the children of $\frac ab$ are $\frac a{a+b}$ (left) and
$\frac{a+b}b$ (right). Number the nodes level by level, left to right: $q_1=\frac11$, $q_2=\frac12$, $q_3=\frac21$, $q_4=\frac13$, \ldots, so that the children of
$q_n$ are $q_{2n}$ and $q_{2n+1}$. (a) Show that every positive rational appears exactly once in the tree, in lowest terms.
(b) Show Newman's formula $q_{n+1}=\dfrac1{2\lfloor q_n\rfloor+1-q_n}$, which lists all positive rationals without ever computing a gcd.
\end{exercise}

\section{Hints}

\begin{hint}{ex:sqrt2}
If $q^2<2$ then $\bigl(\frac{2q+2}{q+2}\bigr)^2-2=\frac{2(q^2-2)}{(q+2)^2}$. A least upper bound $s$ would satisfy $s^2=2$.
\end{hint}
\begin{hint}{ex:algebraic}
Count the polynomials with $\deg P+\sum|a_i|=n$.
\end{hint}
\begin{hint}{ex:cauchyfe}
Additivity gives $f(nx)=nf(x)$; apply it to $x=\frac mn$. Then use the density of $\Q$.
\end{hint}
\begin{hint}{ex:square}
Use $|[0,1]|=|\mathcal P(\N)|$, and a bijection $\mathcal P(\N)\times\mathcal P(\N)\to\mathcal P(\N)$ through even and odd numbers.
\end{hint}
\begin{hint}{ex:hamel}
(c) Find $x_1,x_2$ with $(x_1,f(x_1))$ and $(x_2,f(x_2))$ linearly independent over $\R$, and look at their combinations with rational coefficients.
\end{hint}
\begin{hint}{ex:cuts}
(a) The union of a family of cuts is a cut if it is bounded by a cut. (b) The inverse map sends $\alpha$ to $\sup\alpha$.
\end{hint}
\begin{hint}{ex:krull}
Order the ideals containing $I$ that do not contain $1$.
\end{hint}
\begin{hint}{ex:functions}
A continuous function is determined by its values on $\Q$. For $\R^\R$, compare with $\mathcal P(\R)$ and with $\mathcal P(\R\times\R)$.
\end{hint}
\begin{hint}{ex:calkinwilf}
(a) Walk up the tree: the parent of $\frac ab$ is $\frac a{b-a}$ or $\frac{a-b}b$. (b) Induction on $n$, according to whether $q_n$ is a left child ($q_n<1$) or a
right child.
\end{hint}

\section{Solutions}

\begin{solution}{ex:sqrt2}
\idea A least upper bound of $A$ would be a square root of $2$, and Volume~1 shows that $\sqrt2$ is irrational. To stay inside $\Q$, use an
explicit map that improves both lower and upper approximations, so that no rational can be the best one.
\writeup Let $\phi(q)=\frac{2q+2}{q+2}$ for $q>0$. Then $\phi(q)-q=\frac{2-q^2}{q+2}$ and $\phi(q)^2-2=\frac{2(q^2-2)}{(q+2)^2}$. If $q\in A$, then $\phi(q)>q$ and $\phi(q)^2<2$: $\phi(q)\in A$
is larger. If $q>0$ is a rational upper bound of $A$, then $q^2\ge2$ (if $q^2<2$, $\phi(q)\in A$ exceeds $q$), hence $q^2>2$ (no rational squares to
$2$), and then $\phi(q)<q$ with $\phi(q)^2>2$, so $\phi(q)$ is a smaller upper bound (every $a\in A$ has $a^2<2<\phi(q)^2$). So no element of $\Q$ is a
least upper bound: an upper bound $s\in\Q$ has $s>0$ (as $1\in A$) and $\phi(s)$ is smaller.
\end{solution}

\begin{solution}{ex:algebraic}
\idea A countable union of finite sets is countable. Group the polynomials by a size that allows only finitely many of each size.
\writeup For $P=\sum_{i\le d}a_iX^i$ with $a_d\ne0$, let $h(P)=d+\sum|a_i|\in\N$. For each $n$ there are finitely many $P$ with $h(P)=n$ (degree $\le n$, coefficients in
$[-n,n]$), each with at most $n$ roots. So the algebraic numbers are a countable union of finite sets, hence countable
(Proposition~\ref{prop:countable}(c); no choice is needed here, since each finite set of roots can be listed by increasing real part, then
imaginary part). If the transcendental reals $T$ were countable, $\R=T\cup(\R\cap\text{algebraic})$ would be countable, contradicting
Theorem~\ref{thm:continuum}.
\end{solution}

\begin{solution}{ex:cauchyfe}
\idea Additivity forces linearity over $\Q$; continuity or monotonicity transfers it to $\R$ through rational approximations.
\writeup $f(0)=f(0)+f(0)$ gives $f(0)=0$, and $f(-x)=-f(x)$. By induction $f(nx)=nf(x)$ for $n\in\N$. For $q=\frac mn$, $nf(q)=f(m)=mf(1)$, so $f(q)=qf(1)$. If
$f$ is continuous and $q_k\to x$ are rationals, $f(x)=\lim q_kf(1)=xf(1)$. If $f$ is non-decreasing, for rationals $q<x<q'$ we get
$qf(1)\le f(x)\le q'f(1)$, and letting $q,q'\to x$ gives $f(x)=xf(1)$; non-increasing is the same with reversed inequalities.
\end{solution}

\begin{solution}{ex:square}
\idea Both sets are as large as $\mathcal P(\N)$ or $\mathcal P(\N)\times\mathcal P(\N)$, and two copies of $\N$ fit into one. Schr\"oder--Bernstein turns
injections into a bijection.
\writeup By Theorem~\ref{thm:continuum} there is a bijection $\beta:[0,1]\to\mathcal P(\N)$. The map $(A,B)\mapsto\{2a:a\in A\}\cup\{2b+1:b\in B\}$ is a bijection
$\mathcal P(\N)^2\to\mathcal P(\N)$, since a subset of $\N$ is determined by its even and its odd elements. Composing, $[0,1]^2\to\mathcal P(\N)^2\to\mathcal P(\N)\to[0,1]$ is a
bijection. For $\R^n$: $|\R|=|[0,1]|$, so $|\R^2|=|[0,1]^2|=|[0,1]|=|\R|$, and by induction $|\R^n|=|\R^{n-1}\times\R|=|\R\times\R|=|\R|$. Cantor, who
proved this in 1877, wrote to Dedekind: ``I see it, but I do not believe it.'' Such a bijection cannot be continuous
(Exercise~\ref{ex:notR2}), although a continuous surjection exists (Chapter~4).
\end{solution}

\begin{solution}{ex:hamel}
\idea Zorn gives the basis; a coordinate function is $\Q$-linear but not $\R$-linear. Density comes from the fact that two $\R$-independent
vectors have $\Q$-combinations everywhere in the plane.
\writeup (a) Theorem~\ref{thm:basis} with $K=\Q$, $V=\R$. (b) Every $x$ is uniquely a finite combination $\sum_{b\in B}c_b(x)b$ with rational $c_b(x)$, and
$c_b(x+y)=c_b(x)+c_b(y)$ by uniqueness, so $f=c_{b_0}$ is additive. If $f(x)=cx$ for all $x$, then $c=f(b_0)/b_0=1/b_0\ne0$; but $B$ is infinite
($\R$ is uncountable and a finite-dimensional $\Q$-space is countable), so there is $b_1\in B\setminus\{b_0\}$, and $f(b_1)=0\ne cb_1$.
(c) Let $f$ be additive, so $\Q$-linear by Exercise~\ref{ex:cauchyfe}, and not of the form $cx$. Then $\frac{f(x)}x$ is not constant on $x\ne0$: pick $x_1,x_2$ with
$\frac{f(x_1)}{x_1}\ne\frac{f(x_2)}{x_2}$. The vectors $v_i=(x_i,f(x_i))$ are linearly independent over $\R$ (their determinant is $x_1f(x_2)-x_2f(x_1)\ne0$). So
$(s,t)\mapsto sv_1+tv_2$ is a linear bijection of $\R^2$, continuous with continuous inverse, and maps the dense set $\Q^2$ onto a dense set. But
$q_1v_1+q_2v_2=(q_1x_1+q_2x_2,f(q_1x_1+q_2x_2))$ lies on the graph. Hence the graph is dense, and in particular, over any interval $I$, it
contains points $(x,f(x))$ with $x\in I$ and $f(x)$ as large as we like.
\end{solution}

\begin{solution}{ex:cuts}
\idea A cut is the set of rationals below a real, so the union of cuts is the cut of a supremum, and conversely the supremum of a cut is the
real it describes.
\writeup (a) Let $(\alpha_i)$ be a non-empty family of cuts, all contained in a cut $\gamma$, and $\alpha=\bigcup\alpha_i$. Then $\alpha$ is non-empty, $\alpha\subseteq\gamma\ne\Q$,
closed downwards (each $\alpha_i$ is), and has no largest element (an element of $\alpha$ lies in some $\alpha_i$, which contains a larger one). So $\alpha$ is
a cut, it contains each $\alpha_i$, and every cut containing all $\alpha_i$ contains $\alpha$: it is the least upper bound.
(b) For $x\in\R$, $L(x)=\{q:q<x\}$ is non-empty and not all of $\Q$ (Archimedes), closed downwards, without largest element (density). If
$x<y$, a rational between them lies in $L(y)\setminus L(x)$: the map is strictly increasing, hence injective. Onto: let $\alpha$ be a cut, and let $r\notin\alpha$.
Every $a\in\alpha$ satisfies $a<r$ (otherwise $r\le a$ gives $r\in\alpha$), so $\alpha$ is bounded; let $x=\sup\alpha$. If $q\in\alpha$, some $q'\in\alpha$ exceeds $q$, so
$q<q'\le x$. If $q<x$, $q$ is not an upper bound of $\alpha$, so $q<a$ for some $a\in\alpha$, and $q\in\alpha$. So $\alpha=L(x)$.
\end{solution}

\begin{solution}{ex:krull}
\idea Maximal ideals are maximal elements of a poset, so Zorn applies once chains have upper bounds. The number $1$ is what keeps unions
proper.
\writeup Let $P$ be the set of ideals $J\supseteq I$ with $J\ne R$, that is $1\notin J$, ordered by inclusion; $I\in P$. If $(J_k)$ is a non-empty chain in $P$,
its union $J$ is an ideal (two elements lie in a common $J_k$, so their sum does, and so does any multiple), contains $I$ and not $1$:
it is an upper bound in $P$. By Zorn, $P$ has a maximal element $M$. If $M\subsetneq J\ne R$ for an ideal $J$, then $J\in P$, contradicting maximality: $M$ is a
maximal ideal.
\end{solution}

\begin{solution}{ex:functions}
\idea Continuous functions are cheap: countably many values decide them. All functions are expensive: indicators of subsets of $\R$
are functions.
\writeup Write $\mathfrak c=|\R|=|\mathcal P(\N)|$. Restriction to $\Q$ is injective on $C(\R,\R)$, by density and continuity. So $|C(\R,\R)|\le|\R^\Q|=|\R^\N|$.
Now $|\R^\N|=|\mathcal P(\N)^\N|$, and a sequence of subsets $(A_k)$ of $\N$ is the same as the subset $\{(k,a):a\in A_k\}$ of $\N\times\N$, so
$|\mathcal P(\N)^\N|=|\mathcal P(\N\times\N)|=|\mathcal P(\N)|=\mathfrak c$. Constant functions give $\mathfrak c\le|C(\R,\R)|$, and Schr\"oder--Bernstein concludes. For all functions,
$A\mapsto\mathbf 1_A$ injects $\mathcal P(\R)$ into $\R^\R$, and a function is a subset of $\R\times\R$, so $|\R^\R|\le|\mathcal P(\R\times\R)|=|\mathcal P(\R)|$ by Exercise~\ref{ex:square}.
Finally $|\mathcal P(\R)|>|\R|$ by Cantor's theorem.
\end{solution}

\begin{solution}{ex:calkinwilf}
\idea Going up the tree is Euclid's subtractive algorithm, which ends at $\frac11$; that gives existence and uniqueness at once. For (b), the
successor of a node is its sibling (if it is a left child) or the left child of its parent's successor (if it is a right child); the
formula only needs to be checked in these two cases.
\writeup (a) Since $\gcd(a,a+b)=\gcd(a+b,b)=\gcd(a,b)$, every node is in lowest terms, by induction from $\frac11$. A left child $\frac a{a+b}$ is $<1$, a right
child $\frac{a+b}b$ is $>1$, and the root is $1$. So a node's value determines whether it is the root, a left child or a right child, and
determines its parent's value: $\frac a{b-a}$ if $a<b$, $\frac{a-b}b$ if $a>b$. We prove by induction on $a+b$ that each $\frac ab$ in lowest terms occurs exactly
once. For $a+b=2$ it is $\frac11$, which occurs only at the root (every other node is $\ne1$). For $a+b>2$, $a\ne b$ (lowest terms), the would-be
parent is in lowest terms with a smaller sum, so it occurs exactly once, and $\frac ab$ is its left or right child: $\frac ab$ occurs. Two
occurrences would have the same parent value, hence the same parent node, and be on the same side: they coincide.

(b) Let $s(x)=\frac1{2\lfloor x\rfloor+1-x}$. We prove $q_{n+1}=s(q_n)$ by induction on $n\ge1$. For $n=1$: $s(1)=\frac12=q_2$. Let $n\ge2$. If $n=2m$, then $q_n=\frac a{a+b}<1$
is a left child of $q_m=\frac ab$, and $q_{n+1}=\frac{a+b}b$ is its sibling; indeed $s(q_n)=\frac1{1-q_n}=\frac{a+b}b$. If $n=2m+1$, then $q_n=x+1$ with $x=q_m$, and $q_{n+1}=q_{2(m+1)}$ is
the left child of $q_{m+1}$, which is $\frac y{y+1}$ for $y=q_{m+1}$ (the left child of $\frac cd$ is $\frac c{c+d}=\frac{c/d}{c/d+1}$). By induction ($m+1\le n$),
$y=s(x)=\frac1{2k+1-x}$ with $k=\lfloor x\rfloor$, so $q_{n+1}=\frac1{1+1/y}=\frac1{2k+2-x}$. And $s(q_n)=\frac1{2(k+1)+1-(x+1)}=\frac1{2k+2-x}$. Both cases give $q_{n+1}=s(q_n)$.
\end{solution}

\section*{Sources and further reading}

\begin{itemize}
\item P.~R.~Halmos, \emph{Naive Set Theory}, Van Nostrand, 1960. A hundred pages, from the axioms to Zorn's lemma and the
cardinal numbers; the proof of Zorn above follows its spirit.
\item H.~B.~Enderton, \emph{Elements of Set Theory}, Academic Press, 1977. The construction of $\N$, $\Z$, $\Q$ and $\R$ inside set theory,
with every detail.
\item E.~Landau, \emph{Grundlagen der Analysis}, 1930 (English: \emph{Foundations of Analysis}, Chelsea, 1951). The arithmetic of
$\N$, $\Z$, $\Q$, $\R$ and $\C$ from Peano's axioms, proved line by line.
\item R.~Dedekind, \emph{Stetigkeit und irrationale Zahlen}, 1872, and \emph{Was sind und was sollen die Zahlen?}, 1888. Cuts, and the
recursion theorem.
\item G.~Cantor, \emph{\"Uber eine Eigenschaft des Inbegriffes aller reellen algebraischen Zahlen}, J.~reine angew. Math. 77
(1874), 258--262, and \emph{\"Uber eine elementare Frage der Mannigfaltigkeitslehre}, 1891, where the diagonal argument appears.
\item E.~Zermelo, \emph{Beweis, dass jede Menge wohlgeordnet werden kann}, Math. Annalen 59 (1904), 514--516; M.~Zorn, \emph{A remark
on method in transfinite algebra}, Bull. AMS 41 (1935), 667--670.
\item N.~Calkin and H.~S.~Wilf, \emph{Recounting the rationals}, American Math. Monthly 107 (2000), 360--363. Exercise~\ref{ex:calkinwilf}, with
Newman's formula.
\item T.~Jech, \emph{The Axiom of Choice}, North-Holland, 1973 (Dover, 2008). What choice can and cannot do, and the models of
G\"odel and Cohen.
\item \emph{For the stories in the introduction.} Dedekind dates his discovery to 24 November 1858 in the preface of \emph{Stetigkeit und irrationale Zahlen}. The letters
between Cantor and Dedekind are in E.~Noether and J.~Cavaill\`es (eds.), \emph{Briefwechsel Cantor--Dedekind}, Hermann, 1937; on what ``je le vois, mais je ne le crois
pas'' (letter of 29 June 1877) really meant, see F.~Q.~Gouv\^ea, ``Was Cantor surprised?'', \emph{American Math. Monthly} 118 (2011), 198--209.
\end{itemize}
